% =====================================================================
% RELATORIO - PBL CAFEICULTURA (GRUPO 2)
% Modelagem Estatistica | PPGCA/UNIVALI | Prof. Dr. Rodrigo Sant'Ana
%
% COMPILAR:  Rscript relatorio/build.R
%            (nao compile este arquivo direto: o build regera numeros.tex antes)
%
% ESTRUTURA:
%   main.tex      este arquivo  - preambulo e montagem. Raramente muda.
%   conteudo.tex  o texto       - e aqui que voce escreve.
%   numeros.tex   GERADO        - nao editar, e sobrescrito a cada build.
%   refs.bib      bibliografia
%
% LIMITE: 8 paginas, conforme o enunciado (secao 5).
% =====================================================================
\documentclass[12pt,a4paper]{article}

% --- idioma e codificacao ---------------------------------------------
\usepackage[utf8]{inputenc}
\usepackage[T1]{fontenc}
\usepackage[brazilian]{babel}
\usepackage{lmodern}

% --- layout ABNT: margens 3/2/3/2 cm ----------------------------------
\usepackage[top=3cm,bottom=2cm,left=3cm,right=2cm]{geometry}
\usepackage{indentfirst}          % ABNT indenta o primeiro paragrafo
\usepackage{setspace}
\onehalfspacing                   % entrelinhas 1,5 (ABNT NBR 14724)

% --- matematica, tabelas, figuras -------------------------------------
\usepackage{amsmath,amssymb}
\usepackage{booktabs}
\usepackage{multirow}
\usepackage{graphicx}
\usepackage{float}
\usepackage[font=small,labelfont=bf]{caption}
\graphicspath{{figs/}}

% --- citacao ABNT (NBR 6023), sistema autor-data ----------------------
% Usa biblatex-abnt, nao abntex2cite: o abntex2cite esta sem manutencao e
% quebra com o kernel LaTeX atual (colisao de \abntnextkey com o expl3).
% O backend e o biber, que o TinyTeX instala junto.
\usepackage{csquotes}   % o biblatex recomenda quando ha babel carregado
\usepackage[backend=biber,style=abnt,repeatfields=false]{biblatex}
\addbibresource{refs.bib}

% --- links ------------------------------------------------------------
\usepackage[hidelinks]{hyperref}
\usepackage{url}

% --- numeros gerados pelos modelos ------------------------------------
% Um \newcommand por valor, produzido por relatorio/numeros.R.
% No texto, escreva \bPhi em vez de 20,20. Se o modelo mudar, o texto
% acompanha sozinho. Ver o README para a lista completa de macros.
% =====================================================================
% ARQUIVO GERADO AUTOMATICAMENTE - NAO EDITAR A MAO
% Fonte: relatorio/numeros.R  |  Regerar: Rscript relatorio/numeros.R
% Qualquer edicao aqui e perdida na proxima compilacao.
% =====================================================================

\newcommand{\baseN}{340}
\newcommand{\baseVars}{19}
\newcommand{\baseEsforcoMin}{60}
\newcommand{\baseEsforcoMax}{360}
\newcommand{\baseEsforcoRaz}{6}
\newcommand{\aPearsonRazao}{3,174}
\newcommand{\aDispersaoP}{$<$ 0,001}
\newcommand{\aDispersaoEst}{3,091}
\newcommand{\aGamaEst}{0,898}
\newcommand{\aGamaIcLo}{0,816}
\newcommand{\aGamaIcHi}{0,980}
\newcommand{\aGamaContemUm}{nao contem}
\newcommand{\aTheta}{4,370}
\newcommand{\aThetaSE}{0,542}
\newcommand{\aAIC}{1950,13}
\newcommand{\aAICpois}{2261,25}
\newcommand{\aNpar}{7}
\newcommand{\aPearsonRazaoNB}{1,085}
\newcommand{\aIRRAdubacaon}{0,9992}
\newcommand{\aIRRAltitude}{0,9989}
\newcommand{\aIRRDensidadeplantio}{1,0002}
\newcommand{\aIRRIdadelavoura}{1,0173}
\newcommand{\aIRRMateriaorganica}{1,1404}
\newcommand{\aIRRUmidaderelativa}{1,0498}
\newcommand{\bYmin}{0,0010}
\newcommand{\bYmax}{0,6098}
\newcommand{\bYdistintos}{320}
\newcommand{\bYzeros}{0}
\newcommand{\bYuns}{0}
\newcommand{\bPhi}{20,20}
\newcommand{\bPhiFull}{20,72}
\newcommand{\bPseudoR}{0,4980}
\newcommand{\bLogLik}{455,59}
\newcommand{\bAIC}{-889,18}
\newcommand{\bAICfull}{-879,47}
\newcommand{\bBIC}{-847,06}
\newcommand{\bBICfull}{-802,89}
\newcommand{\bNpar}{10}
\newcommand{\bNparFull}{19}
\newcommand{\bCoefCatuai}{-0,5136}
\newcommand{\bORCatuai}{0,5983}
\newcommand{\bORloCatuai}{0,4987}
\newcommand{\bORhiCatuai}{0,7178}
\newcommand{\bPCatuai}{$<$ 0,001}
\newcommand{\bCoefIcatu}{-1,2805}
\newcommand{\bORIcatu}{0,2779}
\newcommand{\bORloIcatu}{0,2200}
\newcommand{\bORhiIcatu}{0,3510}
\newcommand{\bPIcatu}{$<$ 0,001}
\newcommand{\bCoefMundoNovo}{-0,3093}
\newcommand{\bORMundoNovo}{0,7339}
\newcommand{\bORloMundoNovo}{0,6167}
\newcommand{\bORhiMundoNovo}{0,8734}
\newcommand{\bPMundoNovo}{$<$ 0,001}
\newcommand{\bBourbonPct}{20,80}
\newcommand{\bIcatuPct}{6,82}
\newcommand{\bDifPP}{13,97}
\newcommand{\bLRphi}{3,344}
\newcommand{\bPphi}{0,342}
\newcommand{\bDaicCauchit}{18,17}
\newcommand{\bDaicLoglog}{3,00}
\newcommand{\bDaicLogit}{0,25}
\newcommand{\bLigacoesEquiv}{4}
\newcommand{\bPPligacaoMin}{11,0}
\newcommand{\bPPligacaoMax}{12,9}
\newcommand{\bPPfavoravel}{22,0}
\newcommand{\bPPdesfavoravel}{8,0}
\newcommand{\bRazaoMu}{4,5}
\newcommand{\bRazaoVar}{4,7}
\newcommand{\bRazaoSdResid}{1,210}
\newcommand{\bBartlettP}{0,543}
\newcommand{\bBrownForsytheP}{0,552}
\newcommand{\bPhiIcLo}{0,57}
\newcommand{\bPhiIcHi}{2,31}
\newcommand{\cExporta}{160}
\newcommand{\cNaoExporta}{180}
\newcommand{\cPrevalencia}{47,06}
\newcommand{\cEPV}{16}
\newcommand{\cTotalCatuai}{0,0028}
\newcommand{\cDiretoCatuai}{-0,2543}
\newcommand{\cTotalIcatu}{-0,6080}
\newcommand{\cDiretoIcatu}{-1,1488}
\newcommand{\cTotalMundoNovo}{-0,4559}
\newcommand{\cDiretoMundoNovo}{-0,6403}
\newcommand{\cSevCoef}{-3,6833}
\newcommand{\cSevP}{0,002}
\newcommand{\cStatus}{preliminar}


% --- marcador de pendencia --------------------------------------------
% Destaca em vermelho o que ainda falta. Antes da entrega final, procure
% por \pendente no conteudo.tex e remova as chamadas. Para conferir se
% sobrou alguma:  grep -c pendente relatorio/conteudo.tex
\usepackage{xcolor}
\newcommand{\pendente}[1]{\textcolor{red}{\textbf{[PENDENTE: #1]}}}

% --- atalhos de escrita -----------------------------------------------
\newcommand{\var}[1]{\texttt{#1}}                  % nome de variavel
\newcommand{\pkg}[1]{\textsf{#1}}                  % nome de pacote R
\renewcommand{\arraystretch}{1.15}

% =====================================================================
\begin{document}

% --- capa simplificada (relatorio de disciplina, nao monografia) ------
\begin{center}
    {\large UNIVERSIDADE DO VALE DO ITAJAÍ --- UNIVALI}\\[2pt]
    {\normalsize Escola Politécnica}\\[2pt]
    {\normalsize Programa de Pós-Graduação em Computação Aplicada --- PPGCA}\\[2pt]
    {\normalsize Disciplina: Modelagem Estatística}\\[2pt]
    {\normalsize Prof. Dr. Rodrigo Sant'Ana}\\[18pt]

    {\Large\bfseries Modelos Lineares Generalizados aplicados\\
     ao manejo fitossanitário da cafeicultura}\\[6pt]
    {\large Grupo 2 --- Broca, ferrugem e padrão de exportação}\\[18pt]

    {\normalsize Matheus Piotroski Neis \quad$\cdot$\quad
                 Letícia Zorzi Rama \quad$\cdot$\quad
                 André Ribeiro}\\[14pt]

    {\normalsize Itajaí, Santa Catarina --- Setembro de 2026}
\end{center}

\vspace{6pt}
\hrule
\vspace{10pt}

% =====================================================================
% CONTEUDO DO RELATORIO - e aqui que se escreve
%
% REGRA DE OURO: numero no texto vem de macro, nao digitado.
%   errado:  o parametro de precisao foi 20,20
%   certo:   o parametro de precisao foi \bPhi
% Lista completa das macros: Rscript relatorio/numeros.R && cat numeros.tex
%
% MARCACOES DE PENDENCIA - buscar por "PENDENTE" antes de entregar:
%   \pendente{...}  destaca em amarelo no PDF o que ainda falta
% =====================================================================

\section{Introdução}

A cooperativa precisa decidir onde concentrar o esforço de manejo na próxima safra e
que recomendações fazer aos cooperados quanto a cultivar, densidade de plantio,
adubação e irrigação. O conjunto de dados reúne \baseN{} talhões de café descritos por
\baseVars{} variáveis agronômicas, climáticas e fitossanitárias.

As três perguntas do corpo técnico têm naturezas estatísticas distintas e exigem, cada
uma, uma família diferente de modelo linear generalizado \cite{nelder1972,mccullagh1989}:
a infestação pela broca é uma \textbf{contagem} sob esforço amostral variável; a
severidade da ferrugem é uma \textbf{proporção contínua} sem denominador; e o padrão de
exportação é uma resposta \textbf{binária}. Tratar as três com a mesma ferramenta seria
ignorar o que a natureza de cada dado impõe.

Este relatório apresenta os três modelos, o processo de seleção de variáveis, o
diagnóstico de resíduos e a recomendação de manejo decorrente.

\section{Materiais e métodos}

\subsection{Base de dados}

Os \baseN{} talhões trazem, além das covariáveis agronômicas, três desfechos distintos:
\var{n\_brocas\_capturadas}, \var{severidade\_ferrugem} e \var{padrao\_exportacao}. Essa
estrutura exige cuidado: \textbf{cada desfecho é variável resposta em um desafio e
candidato a preditor nos demais}, e nem toda inclusão é legítima. A questão reaparece
explicitamente na Seção~\ref{sec:mediacao}.

Há ainda duas variáveis que descrevem o \emph{processo de medição}, não o talhão:
\var{n\_armadilhas} e \var{dias\_exposicao}. O produto das duas define o esforço
amostral, que varia de \baseEsforcoMin{} a \baseEsforcoMax{} unidades --- uma amplitude
de \baseEsforcoRaz{} vezes. Comparar contagens brutas entre talhões seria, portanto,
comparar coisas diferentes.

\subsection{Procedimentos comuns}

Todos os ajustes foram feitos em R \cite{rcore2026}, com scripts reprodutíveis e semente
fixada. Três procedimentos são comuns aos três desafios:

\begin{enumerate}
    \item \textbf{Seleção de variáveis} por teste da razão de verossimilhanças, partindo
          do modelo amplo, com exclusão prévia das variáveis irrelevantes por construção
          e verificação de multicolinearidade pelo fator de inflação da variância.
    \item \textbf{Diagnóstico} por resíduos quantílicos aleatorizados \cite{dunn1996},
          envelope simulado \cite{atkinson1981}, verificação de dispersão, identificação
          de pontos influentes e exame da linearidade das covariáveis contínuas na escala
          do preditor linear.
    \item \textbf{Interpretação na escala original} do fenômeno --- brocas por unidade de
          esforço, pontos percentuais de severidade e probabilidade de exportação ---,
          e não apenas na escala do preditor linear.
\end{enumerate}

\section{Desafio A --- infestação pela broca}
\label{sec:desafioA}

\pendente{Seção em consolidação: o modelo final da equipe ainda está sendo revisado.
Os números abaixo já vêm das macros e se atualizam sozinhos.}

\subsection{O esforço amostral e o \textit{offset}}

A contagem de brocas depende de quantas armadilhas foram instaladas e por quantos dias.
Um talhão com esforço seis vezes maior tende a capturar mais insetos sem que a
infestação seja maior. A quantidade de interesse é a \emph{taxa} de captura por unidade
de esforço, $\mu/E$, e não a contagem:

\begin{equation}
    \log\!\left(\frac{\mu_i}{E_i}\right) = \mathbf{x}_i'\boldsymbol{\beta}
    \quad\Longleftrightarrow\quad
    \log(\mu_i) = \log(E_i) + \mathbf{x}_i'\boldsymbol{\beta},
    \label{eq:offset}
\end{equation}

\noindent com $E_i = \var{n\_armadilhas}_i \times \var{dias\_exposicao}_i$. O termo
$\log(E_i)$ entra com coeficiente fixado em~1 --- é um \textit{offset}, não um preditor,
e por isso não consome grau de liberdade.

Essa fixação é uma hipótese testável. Deixando o coeficiente livre, obtém-se
$\hat\gamma = \aGamaEst$, com intervalo de confiança de 95\%
$[\aGamaIcLo;\ \aGamaIcHi]$, que \aGamaContemUm{} o valor~1.

\pendente{O intervalo não contém 1 sob a especificação apenas com covariáveis numéricas.
Com os fatores categóricos no modelo, passa a conter. A equipe precisa decidir qual
especificação reportar e justificar a fixação do \textit{offset} à luz disso --- ou
assumir explicitamente a proporcionalidade como hipótese de trabalho.}

\subsection{Da Poisson à binomial negativa}

A Poisson impõe $\mathrm{Var}[y] = \mathbb{E}[y]$. A razão entre a estatística de Pearson
e os graus de liberdade foi \aPearsonRazao{}, muito acima de~1, e o teste de
\textcite{cameron1990} rejeita a equidispersão ($p = \aDispersaoP$). A variância excede
a média, e ignorar isso subestima os erros-padrão.

Adotou-se então a binomial negativa \cite{hilbe2011}, que acrescenta um parâmetro de
forma $\theta$ e admite $\mathrm{Var}[y] = \mu + \mu^2/\theta$. O ajuste resultou em
$\hat\theta = \aTheta$ (erro-padrão \aThetaSE{}), e a razão de Pearson caiu para
\aPearsonRazaoNB{} --- compatível com~1. O AIC passou de \aAICpois{} na Poisson para
\aAIC{} na binomial negativa.

\subsection{Efeitos estimados}

Com ligação logarítmica e \textit{offset}, $\exp(\beta_j)$ lê-se como \textbf{razão de
taxas} de captura por unidade de esforço. As estimativas do modelo final com \aNpar{}
parâmetros indicam que a umidade relativa (razão de taxas \aIRRUmidaderelativa{} por
ponto percentual) e a matéria orgânica do solo (\aIRRMateriaorganica{}) elevam a
infestação, enquanto a altitude a reduz (\aIRRAltitude{}).

\pendente{Diagnóstico de resíduos e figura do envelope simulado do Desafio A.}

\section{Desafio B --- severidade da ferrugem}
\label{sec:desafioB}

\subsection{Por que não binomial nem regressão linear}

A severidade foi medida como fração de área foliar comprometida por análise de imagem.
\textbf{Não existe denominador}: não se contaram folhas lesionadas de um total de $N$.
Os \baseN{} valores ocupam \bYdistintos{} níveis distintos no intervalo
$[\bYmin;\ \bYmax]$, sem nenhuma observação exata em 0 ou 1
(\bYzeros{} e \bYuns{}, respectivamente).

Forçar a família binomial exigiria \emph{inventar} um $n$. Como a variância binomial é
$\mu(1-\mu)/n$, essa escolha arbitrária propaga-se para toda a inferência: variando o $n$
inventado entre 20 e 2000, o erro-padrão do coeficiente da cultivar varia cerca de dez
vezes, sem que dado algum mude. Correções por quase-verossimilhança ou erros-padrão
robustos eliminam essa dependência, mas ao custo de abandonar a verossimilhança --- e com
ela o teste de razão de verossimilhanças que a análise de dispersão exige.

A regressão linear, por sua vez, produz predições negativas e supõe variância constante,
quando a variância de uma proporção é estruturalmente função da média: não há espaço para
variar perto dos limites do intervalo.

\subsection{Regressão beta}

Adotou-se a regressão beta na parametrização de \textcite{ferrari2004}, implementada
por \textcite{cribari2010}, em que $\mathbb{E}[y] = \mu$ e

\begin{equation}
    \mathrm{Var}[y] = \frac{\mu(1-\mu)}{1+\phi},
    \label{eq:beta}
\end{equation}

\noindent com $\phi$ interpretado como \emph{precisão}: quanto maior, mais concentrada a
distribuição em torno de $\mu$. A variância depende da média por construção, e o problema
estrutural do modelo linear desaparece.

A ligação foi escolhida por critério explícito, e não por ser o padrão do pacote.
Comparadas as seis ligações disponíveis, \bLigacoesEquiv{} ficam dentro de
$\Delta$AIC~$< 2$ e os dados não as distinguem; a diferença estimada entre cultivares
permanece entre \bPPligacaoMin{} e \bPPligacaoMax{} pontos percentuais qualquer que seja
a escolha. A \texttt{cauchit} foi eliminada por ajuste ($\Delta$AIC \bDaicCauchit{}) e a
\texttt{log} por garantir apenas o limite inferior do
intervalo. Entre as restantes, adotou-se a \texttt{logit} por ser a única cujos
coeficientes exponenciados correspondem a razões de chances.

\begin{figure}[H]
    \centering
    \includegraphics[width=0.92\textwidth]{B1_faixa_das_predicoes.png}
    \caption{Faixa de predições das três abordagens contra o domínio válido da
             severidade. O modelo linear invade a região impossível abaixo de zero; a
             regressão beta não pode sair do intervalo, por construção.}
    \label{fig:predicoes}
\end{figure}

\subsection{Dispersão constante}

O enunciado pede que se investigue se $\phi$ é constante. Ajustou-se um modelo com
preditor também para a precisão e comparou-se por razão de verossimilhanças. Nenhuma
covariável, isolada ou em conjunto, rejeitou a constância: para a cultivar,
$\mathrm{LR} = \bLRphi{}$ com $p = \bPphi$.

A sustentação positiva vem do exame direto dos resíduos padronizados: ao longo de cinco
faixas do valor ajustado, a média varia \bRazaoMu{} vezes e a variância bruta
\bRazaoVar{} vezes, enquanto o desvio dos resíduos varia apenas \bRazaoSdResid{} vezes,
sem que os testes de Bartlett ($p = \bBartlettP$) ou de Brown-Forsythe
($p = \bBrownForsytheP$) acusem heterogeneidade. Adotou-se,
portanto, $\phi$ constante, estimado em \bPhi{}.

Registre-se o alcance dessa conclusão: o intervalo de confiança das razões de $\phi$
entre cultivares vai de \bPhiIcLo{} a \bPhiIcHi{}. Descarta-se variação ampla de
precisão, não qualquer variação.

\subsection{Seleção e efeitos}

A seleção \textit{backward} por razão de verossimilhanças reduziu o modelo de
\bNparFull{} para \bNpar{} parâmetros, com melhora simultânea de AIC
(\bAICfull{} $\rightarrow$ \bAIC{}) e BIC (\bBICfull{} $\rightarrow$ \bBIC{}). O modelo
final alcança pseudo-$R^2$ de \bPseudoR{}.

\begin{table}[H]
    \centering
    \caption{Efeito da cultivar sobre a severidade da ferrugem, em razões de chances,
             tomando a Bourbon como referência.}
    \label{tab:cultivar}
    \begin{tabular}{lrrl}
        \toprule
        Cultivar & Coeficiente & Razão de chances & IC 95\% \\
        \midrule
        Catuaí      & \bCoefCatuai     & \bORCatuai     & [\bORloCatuai;\ \bORhiCatuai] \\
        Icatu       & \bCoefIcatu      & \bORIcatu      & [\bORloIcatu;\ \bORhiIcatu] \\
        Mundo Novo  & \bCoefMundoNovo  & \bORMundoNovo  & [\bORloMundoNovo;\ \bORhiMundoNovo] \\
        \bottomrule
    \end{tabular}
\end{table}

Como a ligação \texttt{logit} é não linear, o efeito em pontos percentuais depende do
nível das demais covariáveis. Num talhão típico --- umidade e densidade nas medianas ---,
a severidade esperada é \bBourbonPct\% para a Bourbon contra \bIcatuPct\% para a Icatu,
uma diferença de \bDifPP{} pontos percentuais. Sob condições favoráveis à doença essa
diferença sobe a \bPPfavoravel{} pontos; sob condições desfavoráveis, cai a
\bPPdesfavoravel{}.

\begin{figure}[H]
    \centering
    \includegraphics[width=0.85\textwidth]{B5_meionormal_envelope.png}
    \caption{Gráfico meio-normal com envelope simulado por 99 reajustes do próprio
             modelo. A proporção de pontos fora da banda é compatível com o esperado
             para um envelope de 95\%, sem fugas sistemáticas nas caudas.}
    \label{fig:envelope}
\end{figure}

\section{Desafio C --- padrão de exportação}
\label{sec:desafioC}

\pendente{Seção preliminar. O modelo logístico final ainda não foi ajustado; ver o plano
de implementação. Os números abaixo são descritivos e de mediação, já extraídos da base.}

\subsection{A resposta}

Dos \baseN{} talhões, \cExporta{} atingem o padrão de exportação e \cNaoExporta{} não,
o que corresponde a uma prevalência de \cPrevalencia\%. A resposta é, portanto,
\textbf{balanceada} --- fato que importa para a Seção~\ref{sec:corte}, porque desloca a
discussão do ponto de corte do desbalanceamento para o custo assimétrico dos erros.

A classe minoritária com \cExporta{} casos admite, pela regra de dez eventos por
variável \cite{peduzzi1996}, cerca de \cEPV{} parâmetros. O modelo amplo com a severidade
excede esse limite, o que torna a seleção de variáveis uma necessidade, e não uma
preferência estética.

\subsection{Efeito total e efeito direto da cultivar}
\label{sec:mediacao}

A \var{severidade\_ferrugem} é ela própria um desfecho, influenciada pela cultivar.
Incluí-la como preditor da exportação muda o que o coeficiente da cultivar significa
\cite{pearl2009,vanderweele2016}.

Os dados revelam uma configuração instrutiva: a cultivar Icatu apresenta \emph{a menor}
severidade média e, simultaneamente, \emph{a menor} taxa de exportação. Os dois caminhos
causais operam em direções opostas.

\begin{table}[H]
    \centering
    \caption{Coeficientes da cultivar na predição da exportação, com e sem a severidade
             no modelo. Referência: Bourbon.}
    \label{tab:mediacao}
    \begin{tabular}{lrr}
        \toprule
        Cultivar & Efeito total & Efeito direto \\
                 & (sem severidade) & (com severidade) \\
        \midrule
        Catuaí     & \cTotalCatuai     & \cDiretoCatuai \\
        Icatu      & \cTotalIcatu      & \cDiretoIcatu \\
        Mundo Novo & \cTotalMundoNovo  & \cDiretoMundoNovo \\
        \bottomrule
    \end{tabular}
\end{table}

No modelo que inclui a severidade, o coeficiente desta é \cSevCoef{} ($p = \cSevP$): a
ferrugem de fato reduz a chance de exportação. A Icatu adoece menos, o que a beneficia
por esse caminho; mas seu coeficiente direto é \cDiretoIcatu{}, contra \cTotalIcatu{} no
efeito total. A penalidade direta é parcialmente compensada pelo ganho via menor doença.

\textbf{Qual deles responde à pergunta do produtor?} O produtor escolhe a cultivar
\emph{antes do plantio} e herda os dois caminhos --- o direto e o mediado pela ferrugem.
O \textbf{efeito total} é, portanto, o que responde à sua pergunta. O efeito direto
responde a outra, também legítima: quanto restaria da cultivar se a ferrugem fosse
controlada por outro meio.

Cabe a ressalva: a leitura causal dessa decomposição pressupõe ausência de confundidor
não medido entre ferrugem e exportação --- hipótese que \textbf{não é testável} com estes
dados.

\subsection{Capacidade preditiva e ponto de corte}
\label{sec:corte}

\pendente{Matriz de confusão, curva ROC, AUC com intervalo de DeLong \cite{delong1988},
validação cruzada com seleção dentro do \textit{fold}, calibração e escolha do ponto de
corte por custo assimétrico.}

\section{Decisão e recomendação de manejo}

\pendente{Consolidar após o fechamento dos Desafios A e C. A recomendação deve quantificar
os efeitos na escala original: brocas por unidade de esforço, pontos percentuais de
severidade e probabilidade de exportação.}

A escolha da cultivar é o fator de maior impacto sobre a ferrugem: \bDifPP{} pontos
percentuais de severidade separam a Bourbon da Icatu num talhão típico. Essa vantagem,
contudo, não se traduz automaticamente em exportação --- a Seção~\ref{sec:mediacao}
mostra que a Icatu carrega uma penalidade direta que compensa parte do ganho
fitossanitário. A recomendação de cultivar, portanto, não pode ser derivada apenas do
Desafio B.

\section{Considerações finais}

\pendente{Redigir ao final.}


% --- bibliografia -----------------------------------------------------
\printbibliography[title={Referências}]

\end{document}
